\chapter{仿射簇的不可约性}\label{chap:irreducibility-of-affine-varieties}

\noindent\textbf{第 \thechapter{} 章学习目标}
\begin{itemize}
  \item 举出仿射簇的包含, 交和并的例子, 并用相应的理想描述它们.
  \item 证明每个仿射簇都可以写成某个有限多项式集合的零点集.
  \item 定义仿射簇的不可约性, 并在具体例子中判定不可约性.
  \item 求仿射簇的不可约分解.
  \item 描述把根理想与仿射簇之间的词典限制到素理想或极大理想时会发生什么.
\end{itemize}

研究整数时, 一个关键工具是素因数分解的存在性. 在研究多项式时, 有一个类似的工具, 即分解成不可约因子; 在研究有限 Abel 群时也有类似的工具, 即分解成循环群的直和. 这些情形都说明, 可以通过指出一类数学对象中哪些是不可再分的基本对象, 以及一般对象怎样分解成这些基本组成部分, 来理解这类对象. 在本章中, 我们把这一思想应用于仿射簇, 引入不可约仿射簇的概念, 并说明每个仿射簇怎样唯一地分解成有限多个不可约仿射簇的并, 这些不可约仿射簇就是它的不可约分支.

为实现这一目标, 必须先作一些准备. 首先, 由于我们的目标是证明每个仿射簇都能写成它的不可约分支的并, 所以需要对仿射簇关于包含, 交与并等集合论概念的性质有一般的认识, 这将在第~\ref{sec:inclusions-intersections-unions}~节讨论. 此外, 正如一个整数经过反复分解最终可以分解成素因数, 我们也要确保仿射簇的类似过程不会产生由更小的簇构成的无限嵌套链. 这个条件可转化为 $K[x_{1},\ldots,x_{n}]$ 的一个纯代数性质, 即它是 Noether 环, 这是第~\ref{sec:finite-generation}~节的主题. 作好准备之后, 我们在第~\ref{sec:irreducible-affine-varieties}~节引入不可约性的概念, 并在第~\ref{sec:irreducible-decompositions}~节证明每个仿射簇都能唯一地分解为有限多个不可约仿射簇的并.

\par\smallskip
{\footnotesize
\noindent\textcopyright{} The Author(s) 2025\par
\noindent E. Clader, D. Ross, \emph{Beginning in Algebraic Geometry}, Undergraduate Texts in Mathematics,
\href{https://doi.org/10.1007/978-3-031-88819-9_3}{\nolinkurl{https://doi.org/10.1007/978-3-031-88819-9_3}}.\par}
